\documentclass[10pt,a4paper]{amsart}
\usepackage[margin=19mm]{geometry}
\usepackage{amsmath,amssymb,mathtools}
\usepackage[scr=rsfs,cal=euler]{mathalfa}
\usepackage{xfrac}
\usepackage[unicode,hidelinks,bookmarks=false]{hyperref}
\newcommand{\ie}{i.e.\ }
\newcommand{\cf}{cf.\ }
\newcommand{\resp}{respectively\ }
\newcommand{\Subsection}{\subsection}
\providecommand{\nobreakdash}{\nobreak\hbox{-}}
\setlength{\emergencystretch}{2em}
\multlinegap0pt
\usepackage{amssymb}
\usepackage{mathtools}
\usepackage[shortlabels]{enumitem}
\usepackage{xcolor}
\newtheorem{lemma}{Lemma}
\newtheorem{theorem}{Theorem}
\newcommand{\Z}{\mathbb Z}
\title{Iterated-sumset spectra: The complete exponent law and its rank geometry}
\author{Henry Shin}
\date{Source: arXiv:2609.01690v1 (CC BY 4.0)}
\begin{document}
\maketitle

\noindent\emph{Source and licence.} Iterated-sumset spectra: The complete exponent law and its rank geometry. Version arXiv:2609.01690v1 (CC BY 4.0), distributed by arXiv under the Creative Commons Attribution 4.0 International licence, http://creativecommons.org/licenses/by/4.0/, as recorded in the arXiv licence metadata for this exact version and shown on the abstract page https://arxiv.org/abs/2609.01690v1. Full licence notice: \texttt{licenses/arxiv-2609.01690-CC-BY-4.0.txt}.\par\smallskip

\noindent\textit{Editorial scope.} Counted unit: the article's theorem thm:defecttwodiameter (source0.tex lines 5567--5586). The article's complete original proof of it is delivered with it (source0.tex lines 5588--5749, one \texttt{proof} environment of 162 lines). No mathematics is written, repaired or re-derived: the statement and the proof are the article's own lines, in the article's own notation, and no retained line is substituted or re-flowed.  Retained uncounted as the source-local context the proof needs: lines 1568--1570; lines 5124--5145; lines 5140--5144; lines 5165--5167.  Nothing was omitted from any retained span, so the package carries no omission record.  No retained line contains a citation, so the wrapper carries no bibliography and no bibliography entry.  No retained line contains a figure environment, an image-inclusion command or drawing code, so no image is delivered and the package declares no asset dependency.  Editorial formatting only, in addition: an AMS \texttt{amsart} preamble that restates the article's own declarations and macros used by the retained lines, namely the environments \texttt{lemma}, \texttt{theorem} on separate counters, together with the standard packages those lines need.  The retained environments print in this document as Lemma 1, Theorem 1 in order of appearance, whereas the article prints the same items with the numbering of its own full document.  The complete native source of the article as distributed in the arXiv e-print for this exact version is bundled unchanged as \texttt{sources/209136/source0.tex}.
\begin{equation}\label{eq:Hilbert}
 |hA|=\dim_K(S/I_A)_h.
\end{equation}

\begin{lemma}[Four-point relation-ball bound]
\label{lem:fourpoint-relation-ball}
Let
\[
 A=\{0,a,b,D\},\qquad 0<a<b<D,\qquad \gcd(a,b,D)=1,
\]
and suppose that every nonzero homogeneous relation of $A$ has relation
degree at least $H$, where $H\in\mathbb R$ and $H\geq2$.  Then
\[
 D\geq \binom{H+1}{2}.
\]
If equality holds, then, for the gaps
\[
 g_1=a,\qquad g_2=b-a,\qquad g_3=D-b,
\]
one has
\begin{equation}\label{eq:fourpointcriticalgaps}
 g_1+g_3=H,\qquad
 \gcd(g_1,g_3)=1,\qquad
 g_2=\frac{H(H-1)}2.
\end{equation}
\end{lemma}

\begin{equation}
 g_1+g_3=H,\qquad
 \gcd(g_1,g_3)=1,\qquad
 g_2=\frac{H(H-1)}2.
\end{equation}

\begin{equation}\label{eq:fourpointdeterminant}
 \det\Lambda=[\Z^2:\Lambda]=D.
\end{equation}

\begin{theorem}[Exact diameter at defect two]
\label{thm:defecttwodiameter}
For every $h\geq2$,
\begin{equation}\label{eq:defecttwodiameter}
 \boxed{\qquad \eta_2(h)=1+\binom{h+1}{2}.\qquad}
\end{equation}
More explicitly, if
\[
 D_h=1+\binom{h+1}{2},
\]
then equality is attained by
\[
 \{0,1,D_h-h,D_h\}\qquad(h\ \textup{even})
\]
and by
\[
 \left\{0,\frac{h-1}{2},D_h-\frac{h+1}{2},D_h\right\}
 \qquad(h\ \textup{odd}).
\]
\end{theorem}

\begin{proof}
Normalize a putative witness as
\[
 A=\{0,a,b,D\},\qquad 0<a<b<D,\qquad \gcd(a,b,D)=1.
\]
By \eqref{eq:Hilbert},
\[
 \dim_K(I_A)_h=M_h-|hA|=2.
\]
If $I_A$ contained a nonzero binomial of degree $r<h$, multiplication by
that binomial would inject $S_{h-r}$ into $(I_A)_h$.  This would give
\[
 \dim_K(I_A)_h\geq \dim_K S_{h-r}=M_{h-r}\geq M_1=4,
\]
a contradiction.  Consequently every nonzero homogeneous relation of $A$
has degree at least $h$.

Write the three gaps as
\[
 g_1=a,\qquad g_2=b-a,\qquad g_3=D-b,
\]
and put $C=\binom{h+1}{2}$.  The initial-degree conclusion and
Lemma~\ref{lem:fourpoint-relation-ball}, applied with $H=h$, give
\[
 D\geq C.
\]
If equality held, the gaps would satisfy
\eqref{eq:fourpointcriticalgaps}.  We now exclude that equality.
For $c=(c_1,c_2,c_3)\in\Z^3$ satisfying
$c_1g_1+c_2g_2+c_3g_3=0$, the associated homogeneous relation is
\[
 z(c)=(-c_1,c_1-c_2,c_2-c_3,c_3),
\]
of degree
\begin{equation}\label{eq:defecttwotaildegree}
 \delta(c)=\rho(z(c))=
 \frac{|c_1|+|c_1-c_2|+|c_2-c_3|+|c_3|}{2}.
\end{equation}

Suppose first that $h=2m$.  Reflecting $A$ if necessary, write
$g_1=r\leq m$.  By \eqref{eq:fourpointcriticalgaps},
$\gcd(r,2m)=1$, so $r=2j+1$ is odd.  The vector
\[
 c=(-j,1,-m-j)
\]
is orthogonal to
\[
 (g_1,g_2,g_3)=(r,m(2m-1),2m-r)
\]
and satisfies $\delta(c)=m+r$.  If $r<m$, this is below $h$, a
contradiction.  Hence $r=m$; but then $\gcd(r,h)=m$, so $m=1$.
For $h=2$, the only critical set is $\{0,1,2,3\}$, whose ten degree-two
monomials have only seven distinct images.  Its defect is three, not two.

Suppose now that $h=2m+1$.  After reflection take $g_1=r\leq m$.
The three tail vectors
\[
 \begin{aligned}
 c^{(0)}&=(h-r,0,-r),\\
 c^{(1)}&=(-m,1,-m),\\
 c^{(2)}&=(m+1-r,1,-m-r)=c^{(0)}+c^{(1)}
 \end{aligned}
\]
are orthogonal to $(r,mh,h-r)$, and direct substitution into
\eqref{eq:defecttwotaildegree} gives
\[
 \delta(c^{(0)})=\delta(c^{(1)})=\delta(c^{(2)})=h.
\]
Their degree-$h$ binomials are
\[
 \begin{aligned}
 B_0&=x_1^{h-r}x_2^r-x_0^{h-r}x_3^r,\\
 B_1&=x_0^m x_2^{m+1}-x_1^{m+1}x_3^m,\\
 B_2&=x_1^{m-r}x_2^{m+r+1}
      -x_0^{m+1-r}x_3^{m+r}.
 \end{aligned}
\]
For $1\leq r\leq m$, the six displayed monomials are pairwise distinct.
Thus $B_0,B_1,B_2$ are linearly independent in $(I_A)_h$, contradicting
$\dim_K(I_A)_h=2$.  We have proved the strict lower bound
\begin{equation}\label{eq:defecttwolower}
 D\geq C+1=D_h.
\end{equation}

It remains to verify the asserted witnesses.  If $h=2m$, consider
\[
 A_h=\{0,1,D_h-h,D_h\}
\]
and the two relations
\[
 z=(m,-m-1,m,1-m),\qquad
 w=(m,1-m,-m-1,m).
\]
Both have degree $h$, and the determinant of their middle-coordinate
projections is
\[
 \left|\det
 \begin{pmatrix}
  -m-1&m\\
  1-m&-m-1
 \end{pmatrix}\right|=2m^2+m+1=D_h.
\]
The projected relation lattice has determinant $D_h$ by the same index
calculation as in \eqref{eq:fourpointdeterminant}; hence they form a basis of
the full relation lattice.  For $R=uz+vw$,
\[
 R_0+R_2=hu-v,\qquad R_0+R_3=u+hv.
\]
For any zero-sum vector, the absolute value of every coordinate subsum is
at most its degree.  Hence $\rho(R)\leq h$ implies
\[
 |hu-v|\leq h,\qquad |u+hv|\leq h.
\]
Squaring and adding gives
\[
 (h^2+1)(u^2+v^2)\leq2h^2<2(h^2+1).
\]
Thus $(u,v)$ is zero or a signed coordinate unit vector.

If $h=2m+1$, take
\[
 A_h=\{0,m,D_h-(m+1),D_h\}
\]
and the two relations
\[
 z=(m-1,-m,m+2,-m-1),\qquad
 w=(m+1,-m-1,-m,m).
\]
Since $D_h-(m+1)\equiv1\pmod m$, the set $A_h$ is primitive.  Consequently
the middle-coordinate projection of its full relation lattice has determinant
$D_h$, by the index calculation in \eqref{eq:fourpointdeterminant}.
Again both have degree $h$, while
\[
 \left|\det
 \begin{pmatrix}
  -m&m+2\\
  -m-1&-m
 \end{pmatrix}\right|=2m^2+3m+2=D_h.
\]
Their determinant therefore equals the determinant of the projected lattice,
so they form a basis of the full relation lattice.  For $R=uz+vw$,
\[
 R_0+R_2=hu+v,\qquad R_0+R_3=-2u+hv.
\]
If $\rho(R)\leq h$, put $U=hu+v$ and $V=-2u+hv$.
Then $|U|,|V|\leq h$ and
\[
 u=\frac{hU-V}{h^2+2},\qquad
 v=\frac{2U+hV}{h^2+2}.
\]
Since $h\geq3$, these inequalities give $|u|,|v|<2$.  Each of the four
diagonal choices $(u,v)\in\{-1,1\}^2$ violates either $|U|\leq h$ or
$|V|\leq h$.  Thus once more only zero and the signed coordinate unit
vectors remain.

In either parity, the only nonzero relations of degree at most $h$ are
$\pm z$ and $\pm w$, both of degree exactly $h$.  A primitive relation of
degree $h$ has a unique degree-$h$ binomial, with no room for monomial
padding, and the two binomials here have four distinct monomials.  Therefore
$\dim_K(I_{A_h})_h=2$, so $|hA_h|=M_h-2$.  Together with
\eqref{eq:defecttwolower}, this proves the theorem.
\end{proof}

\end{document}
